\section{Pruebas y resultados}

Este apartado contrasta cada hipótesis específica con mediciones obtenidas sobre
conjuntos de prueba que no participaron del entrenamiento. Se explicitan también
las hipótesis que \textbf{no} pudieron evaluarse: presentar un sistema completo
sin declarar qué partes no fueron medidas equivaldría a afirmar más de lo que la
evidencia sostiene.

\subsection{Condiciones experimentales}

\begin{itemize}
    \item \textbf{Hardware:} GPU NVIDIA RTX 3050 Laptop (4 GB de memoria de video),
        almacenamiento NVMe para los $360.488$ recortes de patente y para los
        conjuntos de atributos.
    \item \textbf{Independencia de los conjuntos:} la partición se realizó por
        sesión de grabación y por bloque horario, nunca por imagen. En
        clasificación de atributos y en detección de placas se verificó
        explícitamente que no exista ninguna sesión compartida entre
        entrenamiento y prueba; el número de sesiones compartidas reportado es
        cero en todos los casos.
    \item \textbf{Métricas:} para clasificación se usa el F1-score macro,
        calculado sobre las clases cuya muestra de prueba alcanza al menos $40$
        ejemplares (\texttt{min\_test\_support}); las clases por debajo de ese
        umbral se informan pero no promedian, porque su intervalo de confianza
        supera los $25$ puntos porcentuales. Para reconocimiento de texto se usa
        la exactitud por carácter ($1 - CER$ promedio) y la coincidencia exacta
        (\textit{Exact Match}). Para detección, precisión, exhaustividad,
        mAP@50 y mAP@50--95.
    \item \textbf{Selección de checkpoint:} en los clasificadores de atributos se
        eligió la época según el F1 macro de validación restringido a clases con
        al menos $20$ ejemplares. Utilizar el macro sobre todas las clases
        presentes dejaba la decisión en manos de clases con tres o cuatro
        imágenes, donde un único acierto desplaza el promedio $1/11 = 0{,}09$.
        Corregir ese criterio elevó el F1 macro de color de $0{,}7937$ a
        $0{,}8124$.
    \item \textbf{Repeticiones:} se ejecutó una corrida por clasificador. La
        varianza medida entre corridas idénticas es de $1$ a $3{,}4$ puntos de
        F1 macro, de modo que las diferencias menores a ese rango no se
        interpretan como mejoras.
\end{itemize}

\subsection{Detección de vehículos (H1.1): no evaluada}

La hipótesis $H_{1.1}$ exige un mAP@50 igual o superior al $80\%$ para el detector
de vehículos. \textbf{Esa medición no se realizó} y no se reporta un valor
estimado.

El motivo es de disponibilidad de datos, no de método: el pool congelado contiene
etiquetas de clasificación (tipo y color por recorte), no cajas delimitadoras
verificadas por un operador. Calcular mAP requiere un conjunto de cajas de
referencia humana, que no existe para los vehículos. El detector YOLO26x opera en
el sistema con sus pesos originales de COCO y sin ajuste fino, por lo que tampoco
hay un conjunto de validación propio que reportar.

La consecuencia se asume explícitamente: el error de la etapa de detección de
vehículos se propaga a todas las etapas posteriores y no está cuantificado en este
trabajo. La línea de trabajo inmediata es anotar cajas de vehículos sobre una
muestra del pool, con el mismo procedimiento de revisión humana utilizado para las
placas (Subsección~\ref{subsec:deteccion_placas}).

\subsection{Clasificación de tipo de vehículo (H1.2)}

El conjunto de prueba contiene $633$ recortes de $5$ clases. La taxonomía excluye
la clase \textit{furgón}, fusionada durante el desarrollo: $98$ recortes
etiquetados como tal se reauditaron individualmente y se reasignaron a
\textit{camioneta} ($23$), \textit{vagoneta} ($20$) y \textit{minibús} ($6$),
descartándose $49$ por no ser identificables.

\begin{table}[H]
    \centering
    \caption{Clasificador de tipo: métricas por clase sobre el conjunto de prueba
    ($n = 633$).}
    \begin{tabular}{lrrrr}
        \toprule
        \textbf{Clase} & \textbf{Precisión} & \textbf{Exhaustividad} & \textbf{F1} & \textbf{n} \\
        \midrule
        automóvil  & 0{,}838 & 0{,}748 & 0{,}791 & 159 \\
        camioneta  & 0{,}880 & 0{,}880 & 0{,}880 &  83 \\
        minibús    & 0{,}949 & 0{,}937 & 0{,}943 &  79 \\
        moto       & 0{,}979 & 1{,}000 & 0{,}989 &  46 \\
        vagoneta   & 0{,}816 & 0{,}868 & 0{,}842 & 266 \\
        \midrule
        \textbf{F1 macro} & \multicolumn{2}{c}{} & \textbf{0{,}8887} & 633 \\
        \textbf{Exactitud} & \multicolumn{2}{c}{} & \textbf{0{,}8578} & 633 \\
        \bottomrule
    \end{tabular}
\end{table}

El error está concentrado en un único par: de $90$ clasificaciones incorrectas,
$62$ ($69\%$) corresponden a la confusión entre \textit{automóvil} y
\textit{vagoneta} ($39$ automóviles predichos como vagoneta, $23$ en sentido
inverso). Es la confusión esperable entre dos siluetas de turismo que, desde una
cámara de tráfico, difieren esencialmente en la altura del techo y el volumen
trasero.

\textbf{Contraste:} $H_{1.2}$ exige F1 macro $\geq 0{,}85$; el valor medido es
$0{,}8887$. \textbf{La hipótesis se respalda.}

\subsection{Clasificación de color predominante (H1.3)}

El conjunto de prueba contiene $551$ recortes distribuidos en $11$ clases, de las
cuales solo seis alcanzan el soporte mínimo de $40$ ejemplares.

\begin{table}[H]
    \centering
    \caption{Clasificador de color sobre el conjunto de prueba. Las clases con
    $n < 40$ se listan solo para dejar constancia de su existencia.}
    \begin{tabular}{lrrrrl}
        \toprule
        \textbf{Clase} & \textbf{Precisión} & \textbf{Exhaustividad} & \textbf{F1} & \textbf{n} & \textbf{Evaluada} \\
        \midrule
        blanco    & 0{,}883 & 0{,}917 & 0{,}900 & 132 & sí \\
        gris      & 0{,}778 & 0{,}575 & 0{,}661 &  73 & sí \\
        plata     & 0{,}824 & 0{,}800 & 0{,}812 & 105 & sí \\
        negro     & 0{,}703 & 0{,}738 & 0{,}720 &  61 & sí \\
        rojo      & 0{,}949 & 0{,}921 & 0{,}935 & 101 & sí \\
        azul      & 0{,}887 & 0{,}810 & 0{,}847 &  58 & sí \\
        \midrule
        verde     & 0{,}500 & 0{,}545 & 0{,}522 &  11 & no \\
        amarillo  & 0{,}000 & 0{,}000 & 0{,}000 &   1 & no \\
        marrón    & 0{,}000 & 0{,}000 & 0{,}000 &   2 & no \\
        beige     & 0{,}083 & 0{,}250 & 0{,}125 &   4 & no \\
        naranja   & 0{,}600 & 1{,}000 & 0{,}750 &   3 & no \\
        \midrule
        \textbf{F1 macro (evaluadas)} & \multicolumn{2}{c}{} & \textbf{0{,}8124} & 530 & \\
        \textbf{F1 macro (las 11 clases)} & \multicolumn{2}{c}{} & \textbf{0{,}5637} & 551 & \\
        \textbf{Exactitud} & \multicolumn{2}{c}{} & \textbf{0{,}8022} & 551 & \\
        \bottomrule
    \end{tabular}
\end{table}

Todo el error se concentra en el bloque acromático. Las confusiones dominantes son
\textit{gris} $\leftrightarrow$ \textit{plata} ($9$ y $5$ casos),
\textit{gris} $\leftrightarrow$ \textit{negro} ($5$ y $4$) y
\textit{plata} $\leftrightarrow$ \textit{blanco} ($12$ y $5$). Separar blanco,
plata, gris y negro en capturas nocturnas es un problema de nivel de brillo, no de
tono, y el clasificador no dispone de una referencia de exposición absoluta: el
mismo vehículo gris claro aparece casi blanco bajo iluminación directa y casi
negro a la sombra.

\textbf{Contraste:} $H_{1.3}$ exige F1 macro $\geq 0{,}85$. El valor medido sobre
las clases evaluables es $0{,}8124$. \textbf{La hipótesis no se respalda} y se
acepta la hipótesis nula $H_{0.3}$ para este componente.

Se deja constancia de una ambigüedad del criterio de evaluación definido en el
perfil del proyecto: la expresión ``F1-score macro'' no especifica el tratamiento
de clases con soporte insuficiente. Bajo la lectura estricta (las once clases) el
valor es $0{,}5637$; bajo la lectura restringida a clases evaluables es $0{,}8124$.
Ambas lecturas quedan por debajo del umbral, por lo que la conclusión no depende
de cómo se resuelva la ambigüedad.

\subsection{Detección de placas de matrícula (H1.4, componente de detección)}
\label{sec:resultados_deteccion}

Sobre el conjunto de prueba de $178$ recortes de vehículo ($140$ con placa
visible) provenientes de $58$ bloques de grabación no vistos durante el
entrenamiento:

\begin{table}[H]
    \centering
    \caption{Detector de placas frente al etiquetador automático, sobre el mismo
    conjunto de referencia humana.}
    \begin{tabular}{lrrrr}
        \toprule
        \textbf{Sistema} & \textbf{Precisión} & \textbf{Exhaustividad} & \textbf{mAP@50} & \textbf{mAP@50--95} \\
        \midrule
        YOLO26n ajustado    & 0{,}9857 & 0{,}9714 & 0{,}9898 & 0{,}6329 \\
        PaddleOCR (localizador) & 0{,}9630 & 0{,}8060 & --- & --- \\
        \bottomrule
    \end{tabular}
\end{table}

\textbf{Contraste:} $H_{1.4}$ exige una exhaustividad $\geq 90\%$ en la detección.
El valor medido es $0{,}9714$. \textbf{El componente de detección se respalda.}

El dato que justificó entrenar un detector propio es la diferencia entre
encontrar la placa y delimitarla correctamente: de las $700$ propuestas
automáticas que sí correspondían a una placa, solo $183$ ($26{,}1\%$) resultaron
utilizables sin corrección humana. El etiquetador localiza, pero su caja recorta
los glifos.

\subsection{Reconocimiento de placas (H1.4, componente de reconocimiento)}
\label{sec:resultados_ocr}

\subsubsection{Sobre recortes de placa aislados}

Conjunto de prueba congelado de $101$ placas reales, sin aumento de datos:

\begin{table}[H]
    \centering
    \caption{Reconocimiento sobre recortes de placa aislados ($n = 101$).}
    \begin{tabular}{lrr}
        \toprule
        \textbf{Modelo} & \textbf{Exactitud por carácter} & \textbf{Exact Match} \\
        \midrule
        CRNN (ResNet18 + BiLSTM + CTC) & 0{,}789 & 0{,}356 \\
        PaddleOCR                      & 0{,}601 & 0{,}297 \\
        \bottomrule
    \end{tabular}
\end{table}

El resultado depende fuertemente de la resolución del recorte, y esa dependencia es
la que fija el techo del módulo:

\begin{table}[H]
    \centering
    \caption{Exactitud por carácter según la altura del recorte, CRNN.
    Mismo modelo en todos los tramos.}
    \begin{tabular}{lrrr}
        \toprule
        \textbf{Altura del recorte} & \textbf{n} & \textbf{Char. accuracy} & \textbf{Exact Match} \\
        \midrule
        $< 20$ px      &  17 & 0{,}605 & 0{,}118 \\
        $20$--$30$ px  &  33 & 0{,}824 & 0{,}333 \\
        $30$--$40$ px  &  22 & 0{,}803 & 0{,}500 \\
        $\geq 40$ px   &  29 & 0{,}847 & 0{,}414 \\
        \bottomrule
    \end{tabular}
\end{table}

Un recorte de $40 \times 13$ píxeles contiene siete caracteres, es decir del orden
de cinco píxeles de ancho por glifo: la información no está en la imagen y ningún
modelo puede reconstruirla. Los diecisiete recortes por debajo de $20$ píxeles de
altura arrastran el promedio global.

\subsubsection{Sobre recortes de vehículo (sistema completo)}

La medición anterior asume un recorte perfecto de la placa. La medición del
sistema ensamblado ---detector propio más CRNN--- se realizó sobre una muestra de
$88$ recortes de vehículo cuyo texto fue verificado manualmente por un operador:

\begin{table}[H]
    \centering
    \caption{Sistema completo sobre $88$ recortes de vehículo con texto verificado.
    Detector y lector reciben exactamente la misma entrada.}
    \begin{tabular}{lrr}
        \toprule
        \textbf{Configuración} & \textbf{Exactitud por carácter} & \textbf{Exact Match} \\
        \midrule
        Detector propio + CRNN (producción) & 0{,}737 & 0{,}261 \\
        Caja humana + CRNN                  & 0{,}799 & 0{,}330 \\
        PaddleOCR de extremo a extremo      & 0{,}653 & 0{,}307 \\
        \bottomrule
    \end{tabular}
\end{table}

El detector localizó la placa en $88$ de $88$ casos de esta muestra
($100\%$). La diferencia entre la segunda y la primera fila ---$6{,}2$ puntos de
exactitud por carácter--- es el costo atribuible al error de localización. El
resto es reconocimiento.

Frente a PaddleOCR, el sistema propio gana en exactitud por carácter ($0{,}737$
contra $0{,}653$) y pierde en coincidencia exacta ($0{,}261$ contra $0{,}307$).
La interpretación es consistente con el comportamiento observado durante todo el
desarrollo: el CRNN emite siempre una cadena plausible y acumula aciertos
parciales, mientras que el motor genérico devuelve vacío cuando no alcanza su
umbral de confianza, o la placa completa cuando la reconoce.

\textbf{Contraste:} $H_{1.4}$ es conjuntiva y exige, simultáneamente,
exhaustividad de detección $\geq 90\%$ y exactitud por carácter $\geq 90\%$. El
primer término se cumple ($0{,}9714$); el segundo no ($0{,}789$ sobre recortes de
placa, $0{,}737$ de extremo a extremo). Adicionalmente, el criterio de evaluación
pide Exact Match $\geq 80\%$, valor que alcanza $0{,}356$. \textbf{La hipótesis
$H_{1.4}$ no se respalda en su totalidad} y se acepta $H_{0.4}$ para el
componente de reconocimiento.

\subsubsection{Nota metodológica sobre la selección de época}

El conjunto de prueba se evaluó al final de cada época durante el entrenamiento
del reconocedor. El valor reportado corresponde a la época $500$, fijada de
antemano con el criterio ``entrenar hasta que la curva se aplane'', y no al máximo
observado en la curva, que fue de $0{,}793$ en una época intermedia. Elegir el
máximo habría introducido un sesgo de selección sobre el propio conjunto de
prueba. Con $n = 101$, el error de muestreo es del orden de $\pm 5$ puntos
porcentuales, por lo que la diferencia entre ambas cifras no es interpretable como
una mejora.

\subsection{Determinación de eventos de ingreso y salida (H1.5): no evaluada}

La hipótesis $H_{1.5}$ requiere exactitud $\geq 90\%$ en la clasificación de
eventos de ingreso y salida. Su medición depende de la simulación de extremo a
extremo con las dos cámaras sincronizadas, que no forma parte del bloque de
trabajo documentado aquí. \textbf{No se reporta un valor.}

\subsection{Contraste de hipótesis}

\begin{table}[H]
    \centering
    \caption{Estado de cada hipótesis específica frente a la evidencia medida.}
    \label{tab:contraste_hipotesis}
    \begin{tabular}{llll}
        \toprule
        \textbf{Hipótesis} & \textbf{Criterio} & \textbf{Medido} & \textbf{Veredicto} \\
        \midrule
        $H_{1.1}$ detección de vehículos & mAP@50 $\geq 0{,}80$ & no medido & sin evaluar \\
        $H_{1.2}$ tipo de vehículo       & F1 macro $\geq 0{,}85$ & $0{,}8887$ & respaldada \\
        $H_{1.3}$ color predominante     & F1 macro $\geq 0{,}85$ & $0{,}8124$ & no respaldada \\
        $H_{1.4}$ placas (detección)     & recall $\geq 0{,}90$ & $0{,}9714$ & respaldada \\
        $H_{1.4}$ placas (reconocimiento) & char acc $\geq 0{,}90$ & $0{,}789$ & no respaldada \\
        $H_{1.4}$ placas (Exact Match)   & EM $\geq 0{,}80$ & $0{,}356$ & no respaldada \\
        $H_{1.5}$ eventos ingreso/salida & exactitud $\geq 0{,}90$ & no medido & sin evaluar \\
        \bottomrule
    \end{tabular}
\end{table}

\subsection{Análisis y limitaciones}

\subsubsection{El techo del reconocimiento es el conjunto de datos}

Dos sistemas de arquitectura y entrenamiento distintos, evaluados sobre el mismo
conjunto, convergen en la misma franja de exactitud por carácter: $0{,}789$ el
CRNN entrenado para este dominio y $0{,}601$ PaddleOCR, que nunca vio estas
cámaras. Aun el sistema con mejor desempeño se queda en $0{,}847$ sobre los
recortes de mayor resolución. Junto con la degradación monótona en función de la
altura del recorte, esto sostiene que la restricción dominante es la calidad y el
volumen del material, no la capacidad del modelo.

Con $496$ placas reales ---de las cuales $395$ se destinan a entrenamiento--- el
ajuste fino dispone de muy pocos ejemplos por carácter. La curva de exactitud
seguía aumentando al finalizar las $500$ épocas, lo que indica que el modelo aún
no había agotado la información disponible; el límite aparece antes por cantidad
de datos que por convergencia.

\subsubsection{Decisiones que se documentan por su efecto medido}

Tres decisiones de implementación modificaron el resultado en un orden de magnitud
mayor que cualquier ajuste de hiperparámetros, y se dejan registradas con su
medición porque el sistema anterior era sustancialmente peor:

\begin{itemize}
    \item \textbf{Normalización de contraste (CLAHE):} $0{,}14 \rightarrow 0{,}676$
        de exactitud por carácter en la corrida de dos etapas. Sin ella, las
        capturas nocturnas quedan fuera de la distribución de entrenamiento.
    \item \textbf{Aumento geométrico de recorte:} $0{,}594 \rightarrow 0{,}737$ de
        exactitud por carácter en el sistema completo. El error de encuadre entre
        el detector que corta en producción y el etiquetador que cortó los datos
        de entrenamiento se absorbía como error de reconocimiento.
    \item \textbf{Selección de época por clases con soporte suficiente:}
        $0{,}7937 \rightarrow 0{,}8124$ de F1 macro en el clasificador de color,
        sin cambiar el modelo ni una sola imagen.
\end{itemize}

\subsubsection{Amenazas a la validez}

\begin{itemize}
    \item \textbf{Una sola corrida por clasificador.} La varianza medida entre
        corridas idénticas es de $1$ a $3{,}4$ puntos de F1 macro; diferencias
        menores a ese rango no se interpretan como mejoras.
    \item \textbf{Muestra pequeña en la medición de extremo a extremo.} Los $88$
        recortes con texto verificado implican un error de muestreo de
        aproximadamente $\pm 5$ puntos en exactitud por carácter. La comparación
        con PaddleOCR en esa condición es indicativa, no concluyente.
    \item \textbf{Selección de la muestra de extremo a extremo.} Los $88$ recortes
        provienen de un conjunto donde el etiquetador automático ya había
        propuesto una caja, por lo que la muestra no es estrictamente aleatoria
        respecto de la dificultad.
    \item \textbf{Un solo recinto y un solo conjunto de cámaras.} Todas las
        sesiones provienen de las mismas cámaras. La partición por bloque horario
        controla la fuga por identidad, pero no la variabilidad entre modelos de
        cámara o condiciones de instalación distintas.
    \item \textbf{Etapas no medidas.} La detección de vehículos y la
        determinación de eventos de ingreso y salida no fueron evaluadas. El
        error de la primera se propaga a todas las etapas posteriores sin
        cuantificar.
    \item \textbf{Clases de color no evaluables.} Cinco de las once clases de color
        tienen menos de $40$ ejemplares en prueba, con intervalos de confianza
        superiores a $25$ puntos. Su desempeño no es interpretable con este
        conjunto.
\end{itemize}
